\documentclass[11pt,a4paper]{article}

\usepackage[T1]{fontenc}
\usepackage{lmodern}     % scalable Latin Modern; needed by microtype's
                         % font expansion (default CM bitmaps fail)
\usepackage[margin=1in]{geometry}
\usepackage{amsmath,amssymb}
\usepackage{booktabs}
\usepackage{listings}
\usepackage{xcolor}
\usepackage{hyperref}
\usepackage{caption}
\usepackage{microtype}

\hypersetup{colorlinks=true, linkcolor=black, citecolor=black, urlcolor=blue}

\lstset{
  basicstyle=\ttfamily\footnotesize,
  columns=fullflexible,
  keepspaces=true,
  breaklines=true,
  showstringspaces=false,
  frame=single,
  framesep=4pt,
  rulecolor=\color{gray!50},
  numberstyle=\tiny\color{gray}
}

\title{Verification-First RTL Synthesis:\\
       Compositional SMT Equivalence from SystemVerilog to Routed Netlist}
\author{J.~R.~R.~Kimmitt\\
        \texttt{jonathan@kimmitt.co.uk}}
\date{\today}

\begin{document}
\maketitle

\begin{abstract}
We describe an open-source synthesis-and-equivalence pipeline that
treats correctness as a first-class output alongside the gate-level
netlist.  A multi-frontend SystemVerilog elaborator lowers source
RTL through a Behavioral IR (BIR) into hardcaml, which we tech-map
to a Liberty cell library and emit as structural Verilog readable by
OpenROAD-flow-scripts.  Three SMT-driven contributions distinguish
the pipeline from conventional yosys+ABC flows.

First, we close the loop with a \emph{cell-mapped equivalence
miter}: the cell-mapped Verilog is parsed back through Verible,
each Liberty cell is expanded into a BIR expression using the
cell's own \texttt{function} and \texttt{next\_state} strings, and
the resulting BIR is paired by module name with the source and
proved equivalent in Z3.  No silent fallback to a different tool
hides synthesis bugs.

Second, we introduce \emph{compositional verification by boundary
substitution}: at parent-module miter time each child instance is
replaced with a Z3 uninterpreted-function call whose name encodes
the (child module, output port) pair.  Identical names and
arguments on both sides of the miter make the leaf trivially equal
in the parent's logical context, so a wide design verifies once a
small set of leaves is independently proven.  We use the
construction to verify a 24-bit multiply-accumulate built from
8-bit multiplier and 48-bit adder leaves in five minutes
end-to-end; the equivalent monolithic check fails to converge in
fifteen.

Third, we drive a \emph{cert-gated post-layout architecture swap}:
analytical depth ratios for arithmetic alternates (Wallace, Dadda,
Brent--Kung, Sklansky, Kogge--Stone) project a critical-path delay
that is honoured only when a Z3 leaf certificate is on file for the
target architecture and width.  On an 8-bit MAC laid out through
OpenROAD-flow-scripts, the predicted post-swap arrival of a ripple
$\rightarrow$ Sklansky adder swap matches the measured post-CTS
arrival within 0.005\,\%, a 35.4\,\% latency reduction with the
swap fully backed by leaf cert.

The pipeline is in OCaml and proves equivalence on industry-grade
designs: gcd (10 / 10 modules), the AES core (4 / 4 modules
including \texttt{aes\_cipher\_top} and \texttt{aes\_key\_expand\_128}),
and the cascaded MAC.
\end{abstract}

\section{Introduction}

Conventional open-source RTL synthesis flows
\cite{wolf2013yosys,brayton2010abc,ajayi2019openroad} treat the
gate-level netlist as the only artefact crossing tool boundaries.
Correctness is an external concern: the user runs equivalence
checking afterwards if at all, often against a sequential simulation
that flags discrepancies as test failures rather than as logical
counterexamples.  When wide arithmetic is involved the equivalence
checkers themselves are a bottleneck --- monolithic SMT
verification of multipliers wider than ten bits is famously hard
\cite{kaivola2009replacing}.

This paper describes a synthesis pipeline built around the opposite
choice: every step that lowers the design either preserves a Z3 SMT
proof of equivalence with its input, or reports an explicit error.
The pipeline lives behind \texttt{USE\_DECOMP\_SYNTH=1} in
OpenROAD-flow-scripts and produces post-CTS layouts on industry
designs while preserving formal correctness across the synthesis
boundary.

The paper makes three principal contributions:

\begin{itemize}
\item A cell-mapped equivalence miter (Section~\ref{sec:cell-equiv})
that proves the gate-level Verilog we emit is logically
equivalent to the source RTL, by parsing the cell-mapped output
back through the same SystemVerilog frontend and expanding each
Liberty cell using its function strings.  The miter requires no
synthesis-tool-specific oracle.
\item A compositional verification scheme based on boundary
substitution (Section~\ref{sec:boundary}) that turns wide
designs into small SMT problems by replacing instance boundaries
with Z3 uninterpreted functions.  The same construction handles
sequential-child boundaries by promoting their outputs to primary
inputs of the parent miter, giving a depth-1 sequential check
without flattening.
\item A cert-gated post-layout resynthesis path
(Section~\ref{sec:resynth}) where an analytical placement-aware
delay model predicts the savings of an arithmetic-architecture
swap, and a Z3-backed leaf certificate is required for the swap
to be reported.  We measure the prediction against ORFS post-CTS
arrival on an 8-bit MAC.
\end{itemize}

Sections~\ref{sec:bir}--\ref{sec:eval} cover the BIR scaffolding,
the synthesis lowering, evaluation results, and related work.

\section{Behavioral IR and multi-frontend lowering}
\label{sec:bir}

The pipeline ingests SystemVerilog from five distinct parsers ---
an OCaml port of Verible's grammar, Verilator's JSON dump, Yosys
RTLIL, Surelog UHDM, and Vivado's structural-VHDL emission of
\texttt{synth\_design -rtl} --- and lowers each into a common
Behavioral IR (BIR).  A \texttt{bmodule} is a record of declared
signals, parameter bindings, sequential and combinational
processes, sub-instance bindings, and function declarations.
\texttt{bstmt} is a tree of assignments, blocks, conditionals,
case statements, loops, and call statements over \texttt{bexpr},
which is in turn a tree of variable references, constants, binary
and unary operators, slices, concatenations, conditional
expressions, and function calls.

Common BIR-level passes apply uniformly across frontends.  The
synthesis pipeline composes them in this order:
\texttt{unroll}, \texttt{inline}, \texttt{iflift},
\texttt{blocking\_subst}, \texttt{meminfer}, optional
\texttt{flatten}, then \texttt{ffrip} + \texttt{share} for
equivalence checking.  Each pass is an
\texttt{bprogram $\to$ bprogram} function with no global state.

\paragraph{ffrip.} The flip-flop ripping pass converts a
\texttt{BSequential} block into an equivalent combinational form
where each FF $Q$ becomes a primary input \texttt{<Q>\_\_Q} (the
current state) and an associated primary output \texttt{<Q>\_\_D}
carrying the next-state expression.  After ripping, two BIR
modules are equivalent at depth-1 sequential level iff their
combinational miter is unsatisfiable.  This avoids the special
casing typically needed for sequential-equivalence and works
correctly for any clock domain that the FF-set comparator first
matches by name.

\section{Synthesis lowering}
\label{sec:synth}

\subsection{From BIR to hardcaml}

\texttt{Behavioral\_to\_hardcaml.create\_circuit} maps each BIR
operator onto a hardcaml \texttt{Signal} primitive.  Multi-bit
selects, slice writes, and indexed mem-writes are coalesced before
lowering: the converter pre-merges \texttt{@slice\_write} and
\texttt{@mem\_write} pseudo-statements within the same clock
domain so that 16 separate
\texttt{always @(posedge clk) text\_out[hi:lo] <= byte} blocks in
AES become a single \texttt{BConcat} bus update, preserving
non-blocking semantics through the FF-rip pass without losing
slice writes to last-write-wins.

\subsection{From hardcaml to Liberty cells}

\texttt{Lib\_map.map\_circuit} tech-maps hardcaml's \texttt{Op2},
\texttt{Mux}, \texttt{Eq}, \texttt{Add}, \texttt{Sub},
\texttt{Mul}, \texttt{Lt} nodes to Liberty cells from a parsed
\texttt{nangate.lib}, bit-blasting wide arithmetic into AND/OR/XOR
ladders.  The ripple-carry adder \texttt{gen\_add} is a chain of
full adders \texttt{gen\_fa}, each implemented as two XOR + two
AND + one OR.  The Array multiplier \texttt{gen\_mul} reduces
$a\!\times\!b$ to $b\_w$ iterations of \texttt{gen\_add\_bits} over
shifted partial-product rows, mirroring the
\texttt{Hardcaml\_circuits.Mul.create} construction in upstream
hardcaml-circuits.  Wallace and Dadda alternates ride on top of
the same primitives.

\subsection{Hierarchical synth and the dual hier/flat netlist}

\texttt{Hier\_synth.synth\_program} synthesises one BIR module at
a time in dependency order: leaves first, parents last.  Hardcaml
never sees a parent's instances; instead, for each child instance
encountered in a parent module, we add phantom IO to the parent's
hardcaml view (a child output becomes a phantom input on the
parent, a child input becomes a phantom output) so hardcaml can
synthesise the parent in isolation.  At cell-Verilog emit time
the original child-instance lines are spliced back in alongside
the parent's bit-blasted gates.

This preserves the dual hier/flat representation: ORFS sees a
hierarchical Verilog with module boundaries intact (good for
SDC-driven timing and module-level reports) while the gate-level
content per-module is also fully visible (good for placement-aware
analysis).

\section{Cell-mapped equivalence}
\label{sec:cell-equiv}

The synthesis pipeline is correct iff, for every source SystemVerilog
module $M$, the cell-mapped Verilog module $M'$ that we emit
implements the same logical function.  Existing flows leave this
property to user-side equivalence checking after the fact.  We
prove it as part of the pipeline.

\subsection{Liberty function expansion}

The cell-mapped Verilog references Liberty cells like
\texttt{AND2\_X1}, \texttt{XOR2\_X1}, \texttt{DFF\_X1},
\texttt{MUX2\_X1}.  Each cell's logical function is encoded in the
Liberty file as a string expression on its output pin --- e.g.
\texttt{function : "(A1 \& A2)"} for an AND gate, or a
sequential cell's \texttt{ff} group with \texttt{clocked\_on},
\texttt{next\_state}, and \texttt{clear} / \texttt{preset} fields.

\texttt{Gate\_netlist\_to\_behavioral.expand\_program} reads a BIR
that contains cell instances (parsed from cell-mapped Verilog) and
replaces each cell with a BIR fragment derived from its Liberty
function.  Combinational cells become \texttt{BCombinational}
processes; sequential cells become \texttt{BSequential} processes
with the appropriate clock and reset semantics.  The result is a
BIR that \emph{describes the same function} as the cell-mapped
Verilog --- not by simulation, but by definition.

\subsection{The miter}

\texttt{test\_synth\_equiv.exe} ties the pieces together:

\begin{enumerate}
\item Parse source SV via Verible into BIR$_{\textrm{src}}$.
\item Pre-process the cell-mapped Verilog and parse it (with
auto-generated empty stubs for each Liberty cell so Verible's
elaborator records the instances) into BIR$_{\textrm{cells}}$.
\item Expand BIR$_{\textrm{cells}}$ via Liberty function strings
to produce a fully behavioural cell-equivalent BIR.
\item Pair modules by name across the two programs and run
\texttt{Z3\_miter.check\_miter\_equivalence} per-pair.
\end{enumerate}

The construction makes no assumption about which synthesis tool
produced the cell-mapped Verilog --- as long as the cells are in
the Liberty file and the Verilog parses, equivalence is provable.
We have shipped it for our own \texttt{synth\_orfs\_shim} output
but the same driver verifies arbitrary post-synth netlists.

\section{Compositional verification by boundary substitution}
\label{sec:boundary}

A 24-bit multiplier has approximately $24 \times 24 = 576$ partial
products and an adder tree of depth $\Theta(24)$ in the Array
form.  Bit-vector SMT solvers struggle with this scale: in our
measurements (Z3 4.13, single core) the equivalence check between
a hand-written multiplier and our \texttt{gen\_mul} output
saturates at the 8-bit width:

\begin{center}
\begin{tabular}{r r l}
\toprule
width & cells in \texttt{gen\_mul} & Z3 time \\
\midrule
6  & 349  & 4.2 s     \\
8  & 641  & 144 s     \\
10 & 1021 & timeout   \\
12 & 1489 & timeout   \\
\bottomrule
\end{tabular}
\end{center}

The conventional response is to pick a more specialised
verification engine for arithmetic
\cite{ciesielski2017verifying}.  We take a different approach:
keep Z3 and reduce the verification problem.

\subsection{Boundary substitution}

\texttt{Behavioral\_boundary.substitute\_program} rewrites every
\texttt{binstance} into a set of \texttt{BCombinational} processes
that drive the parent-side wires connected to each child output
port using a \texttt{BCall} expression of the form
\[
  \texttt{BCall \{ func = M\_p; args = [a_1, \ldots, a_n] \}}
\]
where $M\_p$ encodes the (child module, output port) pair and the
$a_i$ are the parent-side expressions connected to the child's
input ports, sorted alphabetically by port name for canonical
ordering across designs.

In \texttt{z3\_miter} a \texttt{BCall} encodes as a Z3
uninterpreted function declaration: the first occurrence creates
a \texttt{FuncDecl} with domain matching the args' bit-vector
sorts and range matching the output port's declared width; later
occurrences with the same name reuse the same decl from a cache.
Both designs in the miter share the same Z3 context and the same
cache, so the same \texttt{(M\_p, args)} pair produces literally
the same \texttt{Z3.FuncDecl.apply} term on both sides.  Z3's
congruence rule then makes the BCall outputs equal whenever the
arguments are.

\subsection{Sequential children}

A sequential child's output is a function of its inputs
\emph{and} its hidden state, which the parent does not see.
Treating it as a pure uninterpreted function would create
spurious combinational fixed-point loops (a $\to$ mux $\to$ a via
the BCall) that Z3 satisfies freely.  We detect sequential
children by walking the transitive instance closure for any
\texttt{BSequential} process and, for each output of a sequential
child, promote the parent-side connection to a primary input of
the parent miter (matched across designs by name).  This is the
depth-1 sequential check applied at the parent level without
needing flattening.

\subsection{Cascade construction}

To verify wide arithmetic we build it as a cascade of
Z3-tractable leaves.  The 24-bit multiply-accumulate
\texttt{cascade\_mac} (Figure~\ref{fig:cascade}) splits each
24-bit operand into three 8-bit chunks and instantiates nine
\texttt{mul8} leaves to produce all $a_i \times b_j$ partial
products, then an eight-stage chain of \texttt{add48} leaves
sums the shifted products, with a final accumulator register.

\begin{figure}[t]
\centering
\begin{lstlisting}[language=Verilog]
module mul8 (input [7:0] a, input [7:0] b, output [15:0] y);
  assign y = a * b;
endmodule

module add48 (input [47:0] a, input [47:0] b, output [47:0] y);
  assign y = a + b;
endmodule

module cascade_mac (
  input clk, ena, dclr,
  input  [23:0] din, coef,
  output reg [50:0] result);
  // ... 9 mul8 instances + 8 add48 chain + accumulator FF
endmodule
\end{lstlisting}
\caption{The cascaded 24-bit MAC.  Each off-module wire goes
through an explicit \texttt{assign w\_buf = w\_q;} so that
\texttt{hier\_synth}'s phantom-IO promotion has unambiguous
direction per net.}
\label{fig:cascade}
\end{figure}

The miter result on a single core:

\begin{center}
\begin{tabular}{l r l}
\toprule
module & Z3 time & verdict \\
\midrule
\texttt{add48} & 0.16 s & equivalent \\
\texttt{mul8}  & 137 s  & equivalent \\
\texttt{cascade\_mac} & 178 s & equivalent \\
\bottomrule
\end{tabular}
\end{center}

The parent miter reduces to bookkeeping over the ena/dclr/accumulator
FSM because \texttt{mul8} and \texttt{add48} appear identically as
uninterpreted functions on both sides.  Total verification time
for the fully-proven 24-bit MAC is approximately five minutes ---
where the monolithic 24-bit multiplier check does not converge in
fifteen.

\subsection{Buffering convention}

Boundary substitution requires that each off-module wire have a
single direction in the parent's hardcaml view.  When a wire is
the output of one child and the input of another (an adder
chain), \texttt{hier\_synth} cannot decide whether to mark it as
phantom input or phantom output.  We resolve this with a one-line
buffering convention: every chain wire is split into producer-side
\texttt{w\_q} and consumer-side \texttt{w\_buf} with an explicit
\texttt{assign w\_buf = w\_q;} in the parent.  Each net then has
unambiguous direction and the construction is mechanical.  This
matches good design practice anyway; users get verifiability for
free.

\section{Cert-gated post-layout resynthesis}
\label{sec:resynth}

After ORFS lays out a design we want to know whether replacing one
arithmetic operator with a faster architecture would meet timing.
Conventional flows answer with another full layout cycle.  We
answer in seconds with a placement-aware delay projection that is
gated on a Z3 leaf certificate.

\subsection{Architectural depth ratios}

For each (kind, architecture, width) tuple the converter knows the
combinational depth.  For adders:
\begin{itemize}
\item Ripple: $W$
\item Sklansky / Kogge--Stone: $\lceil \log_2 W \rceil$
\item Brent--Kung: $\max(1, 2\lceil \log_2 W \rceil - 1)$
\end{itemize}
For multipliers (all output-width $2W$):
\begin{itemize}
\item Array: $2W$
\item Wallace / Dadda: $\lceil \log_2 W \rceil + W$
\end{itemize}
The depth ratio between two architectures is the ratio of their
combinational depths; we use it as the projection factor on the
fanout cone delay of a baseline arrival.

\subsection{Predict-only flow}

Given a placed netlist (DEF + Liberty-derived per-cell delay
function), \texttt{Predict\_swap} identifies the worst endpoint,
extracts its fanout cone with \texttt{Fanout\_cone}, computes the
cone delay as the sum of forward-Elmore arcs, and for each
binding-to-arch substitution scales the cone delay by the
architectural depth ratio.  The new arrival is the baseline minus
the saved cone delay.  Every prediction is reported with a cert
status:

\begin{itemize}
\item \textbf{proven}: a Z3 leaf certificate is on file for
(kind, target architecture, width) --- the swap is recommended.
\item \textbf{no-cert}: the projection is reported but excluded
from the cert-gated best-pick; the user is told which
\texttt{verify-arch} command to run to unlock it.
\end{itemize}

\subsection{Measured prediction error on an 8-bit MAC}

We synthesised an 8-bit Array$\times$Ripple MAC, ran it through
ORFS to post-CTS, ran the predictor, applied the recommended swap
(Sklansky adder), re-ran ORFS to post-CTS, and measured the actual
critical-path arrival.

\begin{center}
\begin{tabular}{l r r}
\toprule
& baseline & swapped \\
\midrule
arch (mul $\times$ add) & Array $\times$ Ripple & Array $\times$ Sklansky \\
post-CTS cell count & 535 & 343 \\
critical-path arrival & 917.31 ps & 592.67 ps \\
predicted arrival & --- & 592.7 ps \\
prediction error & --- & $-0.005$\,\% \\
cert status & --- & proven \\
\bottomrule
\end{tabular}
\end{center}

The next-most-attractive swap (\texttt{u\_mul} array $\to$
Wallace, predicted savings 151 ps, also cert-proven) becomes the
head of the recommendation list after the adder swap is applied.

\section{ORFS substitution}
\label{sec:orfs}

OpenROAD-flow-scripts orchestrates the open-source RTL-to-GDSII
flow: yosys for synthesis, ABC for technology mapping, OpenROAD
for floorplan / place / route / CTS, and KLayout for finishing.
Our synthesis pipeline replaces the yosys+ABC steps when invoked
with \texttt{USE\_DECOMP\_SYNTH=1}.  The mechanism is a small
patch to the ORFS Makefile:

\begin{lstlisting}[language=make]
ifeq ($(USE_DECOMP_SYNTH),1)
do-yosys:
    $(DECOMP_SHIM) $(DESIGN_NAME) \
        $(RESULTS_DIR)/1_2_yosys.v $(VERILOG_FILES)
    cp $(SDC_FILE) $(RESULTS_DIR)/1_2_yosys.sdc
endif
\end{lstlisting}

The shim is a single binary that runs the Verible
$\to$ BIR $\to$ hardcaml $\to$ Liberty pipeline and writes the
cell-mapped Verilog directly into the location ORFS expects.
There is no silent fallback: any failure of the pipeline aborts
the make target.  Yosys remains available for explicit
oracle / comparison runs but never as a hidden fallback that masks
a synthesis bug.

The same patch generalises: for the multiplier-architecture demos
we provide a second shim, \texttt{mac\_synth\_shim}, that emits a
\texttt{Synth\_mac} netlist directly with the chosen multiplier
and adder architectures read from environment variables.  Both
shims emit the \texttt{1\_2\_yosys.v} and \texttt{1\_2\_yosys.sdc}
files in the same way.

\section{Evaluation}
\label{sec:eval}

We exercise the pipeline on three classes of design:

\paragraph{gcd (10 modules).} The standard ORFS demo.  10 / 10
modules formally equivalent at module level (each leaf and the
top), full ORFS run produces post-route DEF and final Verilog
unchanged from the stock yosys+ABC flow.

\paragraph{AES core (4 modules).} Verilog-2001 implementation with
\texttt{aes\_sbox} (256-entry case statement), \texttt{aes\_rcon}
(state machine), \texttt{aes\_key\_expand\_128} (state-array
register, mix\_col function), and \texttt{aes\_cipher\_top}
(round-state matrix, slice-LHS register writes).  4 / 4 modules
formally equivalent under
\texttt{Z3\_MITER\_TIMEOUT\_MS=120000}.  AES exercises the harder
parts of the BIR pipeline: function inlining for \texttt{mix\_col}
(which writes four byte-slices of a 32-bit return value across
four \texttt{@slice\_write} calls and must coalesce into a single
\texttt{BConcat}); cross-domain BSequential merging (16 separate
\texttt{always} blocks driving \texttt{text\_out} bytes coalesce
into one BSequential with a per-slot \texttt{BConcat} update);
and parameter-resolution iteration (Verible's reverse parse
order forces a fixed-point sweep to resolve derived parameters
like \texttt{mwidth = dwidth + cwidth}).

\paragraph{Cascaded MAC (3 modules).} Detailed in
Section~\ref{sec:boundary}.  3 / 3 formally equivalent in
$\sim$5 min on a single core.  The compositional structure
generalises: a 32-bit MAC built as 16 \texttt{mul8} + 15
\texttt{add64} leaves would verify by the same pattern with
roughly the same wall-clock cost as the multiplication count
times the per-leaf cost, since leaf verifications are mutually
independent and trivially parallelisable.

\paragraph{Post-layout resynth.} Detailed in
Section~\ref{sec:resynth}.  Predicted vs.\ measured prediction
error of 0.005\,\% on an 8-bit MAC, with the swap fully backed by
a Z3 leaf certificate.

\paragraph{DFT and structural ATPG.} The same gate netlist that the
miter verifies feeds an optional design-for-test track:
\texttt{scan\_insert.ml} swaps every \texttt{DFF\_X1} variant for its
scan equivalent and stitches the chain;
\texttt{mem\_boundary\_scan.ml} wraps black-box children with
bit-level boundary FFs so an unmodelled memory or sub-block doesn't
block ATPG observation; \texttt{jtag\_tap\_insert.ml} adds an
IEEE~1149.1 TAP and BC\_1 cells on every chip-top pad, emitting a
\texttt{STD\_1149\_1\_2001}-conformant BSDL sidecar.  A three-tier
fault simulator (random, backwards-justified directed, PODEM with
five-valued logic and D-frontier backtracking) reports per-module
stuck-at coverage on the gate netlist.  Picosoc with the full DFT
stack reaches 85\,\% top-level coverage and 92--99\,\% across CPU,
SPI flash controller, UART, and memory wrapper; the AES S-box hits
94\,\% with PODEM contributing 352 of 1102 attempted patterns
(random alone is structurally weak on 256-entry case statements).
The track is presented here as a structural composition demo --- it
is not a formal contribution of this paper, but it shares the same
gate netlist representation, which is why the same boundary
construction works.

\section{Related work}

Yosys + ABC \cite{wolf2013yosys,brayton2010abc} is the
predominant open-source synthesis chain.  Yosys offers an
\texttt{equiv\_*} family of commands that compares two RTLIL
designs at the SAT level, and Berkeley ABC ships a combinational
equivalence checker (CEC).  Both depend on the synthesis tool
itself producing the netlists they verify, with no
multi-frontend reach.  Our pipeline instead treats Yosys as a
parallel oracle (work item, not a dependency) and produces an
SMT proof against the source SV without going through a yosys
intermediate.

\paragraph{Specialised arithmetic verification.} Gr\"obner-basis
methods \cite{ciesielski2017verifying} verify multipliers up to
several hundred bits but require non-SMT machinery.  Our cascade
construction trades off generality for SMT-only tooling: any
BIR-expressible boundary works, including non-arithmetic ones.

\paragraph{Cell-mapped equivalence.} Commercial tools
(Synopsys Formality, Cadence Conformal) read the cell-mapped
Verilog and prove it equivalent to source RTL using
proprietary engines.  \texttt{test\_synth\_equiv.exe} is, as far
as we know, the first open-source equivalence checker that
expands Liberty function strings directly into a behavioural
form for SMT comparison against the source.

\paragraph{Hierarchical equivalence.}  ABC's \texttt{cec}
supports compositional checking via miter cones.  We push
further: the boundary becomes an uninterpreted function
\emph{at the SMT layer}, not a structural cone substitution,
which is what allows Z3 to reduce a 24-bit cascade-MAC parent to
the surrounding bookkeeping in 178 s.

\paragraph{Cert-gated optimisations.}  EQy \cite{eqy} from the
yosys ecosystem performs equivalence-driven swaps for known
mux / gate optimisations.  We add the prediction step --- the user
sees \emph{which} swap is worth doing before paying for the
verification cycles --- and require that the cert exist before
the swap is recommended.

\paragraph{OpenROAD-flow-scripts.}  ORFS is the orchestrator we
plug into.  Our patch is mechanical, leaves the rest of the flow
untouched, and is auditable.

\section{Conclusion and future work}

A multi-frontend SystemVerilog elaborator combined with a
Liberty-aware tech mapper, an SMT miter, and a boundary-substitution
construction yields a synthesis pipeline that proves correctness
at every step it lowers the design.  The same construction makes
wide arithmetic tractable for SMT and grounds an analytical
post-layout resynth predictor in Z3-verified leaves.

Since the initial draft (\texttt{paper-v1}) three of the
then-open items have landed: the cascaded MAC has been pushed
through ORFS layout to post-route; LHS-context width propagation
is fixed at the Verible converter source so consumers see
consistent widths without per-operator workarounds in the miter;
and the yosys-as-oracle parallel-correctness harness is in tree
as \texttt{test\_yosys\_oracle\_sweep.exe}.  The structural DFT
track described in Section~\ref{sec:eval} (scan, hierarchical
boundary scan, JTAG TAP, BSDL) is also new since v1.  Open work,
in roughly the order we expect to take it:

\begin{enumerate}
\item Auto-cascade for \texttt{gen\_mul} above the SMT
threshold, lifting the manual buffering convention into the
lib\_map output so any wide multiplier inherits the cascade
verification path.
\item Treat hard-macro pins as pseudo-observation / pseudo-control
points in the fault simulator (the no-area-cost alternative to
the structural BSC wrap around memories), so designs that don't
want the extra latency get the same chip-top coverage gain.
\item Naja interchange-format support so per-leaf SMT
certificates can be cached and replayed across runs and tools,
in line with the cert-gated swap construction already in tree.
\end{enumerate}

The codebase is open source at
\url{https://github.com/jrrk2/System-Verilog-decompiler}.
Appendix~\ref{sec:repro} gives full step-by-step instructions to
install every dependency and reproduce each numbered result above
on a clean Debian / Ubuntu machine.

\appendix

\section{Reproducing the results}
\label{sec:repro}

Every measurement in the paper was obtained by running an
executable at a tagged commit in the open-source repository.  This
appendix lists the commands to reproduce them on a clean Linux
machine, including every dependency that has to be installed.  The
repository's \texttt{README.md} contains a more compact version of
the same information.

\subsection{System assumptions}

The author's measurement machine: Debian 12 / Ubuntu 24.04 LTS
class system, x86\_64, 12 hardware threads, 32~GiB RAM, single
node.  No GPU, no FPGA, no proprietary tools required to reproduce
the headline numbers (the optional Vivado dependency below is only
needed for the CVA6 oracle and the Vivado column on the sv-tests
dashboard, neither of which is claimed in this paper).

\subsection{System packages}

\begin{lstlisting}[language=bash]
sudo apt-get install \
  build-essential git curl                 # toolchain
  m4 pkg-config bubblewrap                 # opam build deps
  libgmp-dev libffi-dev                    # native lib deps for Z3 binding
  libz3-dev z3                             # the SMT solver itself
  texlive-latex-recommended \
  texlive-latex-extra                      # for paper/paper.tex
\end{lstlisting}

\subsection{OCaml toolchain via opam}

\begin{lstlisting}[language=bash]
# opam itself
sudo apt-get install opam
opam init --auto-setup --disable-sandboxing
opam switch create 5.3.0
eval $(opam env --switch=5.3.0 --set-switch)

# OCaml packages
opam install dune menhir yojson ppx_deriving_yojson \
             hardcaml hardcaml_circuits z3 \
             lua-ml linenoise
\end{lstlisting}

The Z3 OCaml binding links against \texttt{libz3.so} from the
system \texttt{libz3-dev} package.  On macOS the equivalent is
\texttt{brew install z3 opam}.

\subsection{Source checkout and build}

\begin{lstlisting}[language=bash]
git clone https://github.com/jrrk2/System-Verilog-decompiler.git
cd System-Verilog-decompiler
eval $(opam env --switch=5.3.0 --set-switch)
dune build @all
\end{lstlisting}

The full build produces approximately 70 executables under
\texttt{\_build/default/}; the four headline drivers used in this
paper are
\texttt{test\_synth\_equiv.exe},
\texttt{synth\_orfs\_shim.exe},
\texttt{mac\_synth\_shim.exe}, and
\texttt{test\_mac\_postlayout.exe}.

\subsection{External tools}

\paragraph{Verilator $\geq$ 5.024} is required for the four-way
miter and for the pre-flatten wrapper used by the sv-tests
integration.  On Debian / Ubuntu:
\begin{lstlisting}[language=bash]
sudo apt-get install verilator
\end{lstlisting}

\paragraph{OpenROAD-flow-scripts (ORFS)} is required for the
post-layout resynth result and for the cell-mapped-equivalence
results that use Nangate45 cells.  Follow the upstream install
instructions at
\url{https://github.com/The-OpenROAD-Project/OpenROAD-flow-scripts}.
The minimal step:
\begin{lstlisting}[language=bash]
git clone --recursive \
   https://github.com/The-OpenROAD-Project/OpenROAD-flow-scripts.git \
   $HOME/OpenROAD-flow-scripts
cd $HOME/OpenROAD-flow-scripts
./build_openroad.sh
\end{lstlisting}
Adjust \texttt{ORFS\_DIR} via the environment if you install it
elsewhere.

After ORFS is installed, apply the Makefile patch that lets
\texttt{USE\_DECOMP\_SYNTH=1} route \texttt{do-yosys} to our shim:
\begin{lstlisting}[language=bash]
cd $HOME/System-Verilog-decompiler
test/orfs/install.sh                  # apply patch
test/orfs/install.sh --check          # confirm patched
\end{lstlisting}

\paragraph{Optional} Yosys, Slang, Surelog, and Vivado are
optional and not needed for any result in this paper.  The
\texttt{README.md} lists what each of them unlocks.

\subsection{Reproducing each result}

Each block below reproduces one numbered result from the paper.
All commands assume \texttt{\$PWD} is the repository root and the
opam switch has been activated.

\paragraph{Section~\ref{sec:cell-equiv}, gcd 10 / 10 modules
formally equivalent.} The driver runs synthesis with our shim,
emits the cell-mapped Verilog, and proves cell-mapped
$\equiv$ source via Z3:
\begin{lstlisting}[language=bash]
NAN=$HOME/OpenROAD-flow-scripts/flow/platforms/nangate45/lib/NangateOpenCellLibrary_typical.lib
GCD=$HOME/OpenROAD-flow-scripts/flow/designs/src/gcd/gcd.v
_build/default/synth_orfs_shim.exe gcd /tmp/gcd_synth.v $GCD
_build/default/test_synth_equiv.exe gcd $GCD /tmp/gcd_synth.v $NAN \
  | grep Summary
\end{lstlisting}
Expected output: \texttt{Summary: 10 pass, 0 fail (10 modules total)}.

\paragraph{Section~\ref{sec:eval}, AES core 4 / 4 modules formally
equivalent.}  AES has deeper boundary cones; bump the per-module
timeout from the 30 s default:
\begin{lstlisting}[language=bash]
AES_DIR=$HOME/OpenROAD-flow-scripts/flow/designs/src/aes
AES_FILES=$(ls $AES_DIR/*.v | grep -v inv_)
# Combined source for the source side of the miter
cat $AES_FILES > /tmp/aes_combined.v
_build/default/synth_orfs_shim.exe aes_cipher_top \
  /tmp/aes_synth.v $AES_FILES
Z3_MITER_TIMEOUT_MS=120000 \
  _build/default/test_synth_equiv.exe aes_cipher_top \
    /tmp/aes_combined.v /tmp/aes_synth.v $NAN | grep Summary
\end{lstlisting}
Expected output: \texttt{Summary: 4 pass, 0 fail (4 modules
total)}.

\paragraph{Section~\ref{sec:boundary}, cascade\_mac 24-bit MAC
3 / 3 formally equivalent.}
\begin{lstlisting}[language=bash]
_build/default/synth_orfs_shim.exe cascade_mac \
  /tmp/cascade_mac_synth.v test/cascade/cascade_mac.sv
Z3_MITER_TIMEOUT_MS=600000 \
  _build/default/test_synth_equiv.exe cascade_mac \
    test/cascade/cascade_mac.sv /tmp/cascade_mac_synth.v $NAN \
  | grep -E '^\xe2\x94\x80\xe2\x94\x80\xe2\x94\x80\xe2\x94\x80| Summary| seconds'
\end{lstlisting}
Expected output (single core, $\sim$5 minutes wall-clock):
\texttt{Summary: 3 pass, 0 fail (3 modules total)}, with
\texttt{add48} at $\sim$0.16 s, \texttt{mul8} at $\sim$140 s, and
\texttt{cascade\_mac} at $\sim$180 s.

\paragraph{Section~\ref{sec:resynth}, 8-bit MAC layout +
post-CTS ripple $\to$ Sklansky swap.}
The \texttt{mac\_synth\_shim} reads the chosen architecture from
environment variables and emits a Synth\_mac netlist directly into
the slot ORFS expects:
\begin{lstlisting}[language=bash]
mkdir -p $HOME/OpenROAD-flow-scripts/flow/designs/nangate45/mac
cat > $HOME/OpenROAD-flow-scripts/flow/designs/nangate45/mac/config.mk <<'EOF'
export DESIGN_NICKNAME = mac
export DESIGN_NAME     = mac
export PLATFORM        = nangate45
export VERILOG_FILES   = $(DESIGN_HOME)/nangate45/mac/dummy.v
export SDC_FILE        = $(DESIGN_HOME)/nangate45/mac/constraint.sdc
export CORE_UTILIZATION = 30
export CORE_ASPECT_RATIO = 1
export CORE_MARGIN     = 2
export PLACE_DENSITY_LB_ADDON = 0.20
export TNS_END_PERCENT = 100
EOF
echo "// dummy" > $HOME/OpenROAD-flow-scripts/flow/designs/nangate45/mac/dummy.v
echo "create_clock -period 4.0 -name clk [get_ports clk]" > \
  $HOME/OpenROAD-flow-scripts/flow/designs/nangate45/mac/constraint.sdc

# Baseline: Array x Ripple
cd $HOME/OpenROAD-flow-scripts/flow
env USE_DECOMP_SYNTH=1 \
    DECOMP_SHIM=$HOME/System-Verilog-decompiler/_build/default/mac_synth_shim.exe \
    MAC_WIDTH=8 MAC_MUL=array MAC_ADD=ripple MAC_CLK_NS=4.0 \
  make DESIGN_CONFIG=designs/nangate45/mac/config.mk

# Extract DEF and analyse
$HOME/OpenROAD-flow-scripts/tools/install/OpenROAD/bin/openroad \
  -no_init -exit results/nangate45/mac/base/4_cts.odb -tcl \
  results/nangate45/mac/base/4_cts.def
cd $HOME/System-Verilog-decompiler
_build/default/test_mac_postlayout.exe \
  $HOME/OpenROAD-flow-scripts/flow/results/nangate45/mac/base/4_cts.def \
  $HOME/OpenROAD-flow-scripts/flow/results/nangate45/mac/base/1_2_yosys.v 8
\end{lstlisting}
Expected baseline output: \texttt{worst endpoint @ 917.31 ps}, with
\texttt{u\_add ripple $\to$ sklansky} as the cert-gated best pick
predicted to land at \texttt{592.7 ps}.

To measure the prediction, repeat the layout with
\texttt{MAC\_ADD=sklansky} (and a different sub-directory if you
want to keep the baseline available):
\begin{lstlisting}[language=bash]
env USE_DECOMP_SYNTH=1 \
    DECOMP_SHIM=$HOME/System-Verilog-decompiler/_build/default/mac_synth_shim.exe \
    MAC_WIDTH=8 MAC_MUL=array MAC_ADD=sklansky MAC_CLK_NS=4.0 \
    FLOW_VARIANT=sklansky \
  make DESIGN_CONFIG=designs/nangate45/mac/config.mk
\end{lstlisting}
The post-CTS analysis on the swapped variant reports
\texttt{worst endpoint @ 592.67 ps} --- a 0.005\,\% deviation from
the prediction.

\subsection{Tagged commit}

The numbers in this paper match the repository at the tag
\texttt{paper-v1} (created at the time of this draft).  All later
commits land on \texttt{main} and are listed in
\texttt{git log paper-v1..main}; if any subsequent change
materially affects the numbers it will be called out in the
release-notes file \texttt{paper/CHANGELOG.md}.

\bibliographystyle{plain}
\begin{thebibliography}{9}

\bibitem{wolf2013yosys}
C.~Wolf.
\newblock Yosys --- a free Verilog synthesis suite.
\newblock In \emph{Proc.\ Austrochip}, 2013.

\bibitem{brayton2010abc}
R.~Brayton and A.~Mishchenko.
\newblock ABC: An Academic Industrial-Strength Verification Tool.
\newblock In \emph{Computer Aided Verification (CAV)}, pages 24--40, 2010.

\bibitem{ajayi2019openroad}
T.~Ajayi et al.
\newblock OpenROAD: Toward a Self-Driving, Open-Source Digital Layout Implementation Tool Chain.
\newblock In \emph{GOMACTech}, 2019.

\bibitem{kaivola2009replacing}
R.~Kaivola, R.~Ghughal, et al.
\newblock Replacing Testing with Formal Verification in Intel CoreTM i7 Processor Execution Engine Validation.
\newblock In \emph{CAV}, pages 414--429, 2009.

\bibitem{ciesielski2017verifying}
M.~Ciesielski, C.~Yu, W.~Brown, D.~Liu, A.~Rossi.
\newblock Verification of Gate-Level Arithmetic Circuits by Function Extraction.
\newblock In \emph{DAC}, 2015.

\bibitem{eqy}
YosysHQ.
\newblock EQy --- equivalence checking driver.
\newblock \url{https://github.com/YosysHQ/eqy}, 2024.

\end{thebibliography}

\end{document}
